\ifx\danslelivre\undefined
\RequirePackage{luatex85}
\documentclass[10pt]{../fichiers-configuration-latex/smfart}
\usepackage{polyglossia}
%\usepackage[francais]{babel}
\usepackage{fontspec}
\usepackage[math-style=ISO,bold-style=upright,sans-style=upright]{unicode-math}
%\usepackage[mathletters]{ucs}
%\usepackage[utf8]{inputenc}
%\usepackage[T1]{fontenc}
\usepackage[unicode,colorlinks=true,linkcolor=blue,citecolor=Sepia,urlcolor=magenta]{hyperref}
\setmainlanguage{french}
\setmainfont[Ligatures=TeX]{Linux Libertine O}
% variante possible : TeXGyrePagella ?
\setmathfont{XITS Math}

% maths
\usepackage{amsmath}
%\usepackage{amssymb}
%\usepackage{mathrsfs}
%\usepackage{stmaryrd}
%\usepackage{euscript}

% couleurs
\usepackage[usenames,dvipsnames]{xcolor}

% graphiques
\usepackage{graphicx}
\usepackage[all]{xy}
% Xy-pic 3.8.9's PDF driver assumes that \fontdimen8\textfont3 is a
% legacy TeX rule width.  Under LuaTeX with unicode-math/XITS Math it is
% instead 6.62pt, which turns every arrow shaft into a thick black bar.
% Override only xypdf's private cached width; do not alter the math font.
\makeatletter
\AtBeginDocument{%
  \@ifundefined{xP@preclw}{}{%
    \edef\xP@preclw{0.4pt}%
    \edef\xP@lw{\xP@dim\xP@preclw}%
    \PackageInfo{sga4-xypdf-compat}{Xy-pic PDF line width set to \xP@preclw}%
  }%
}
\makeatother
%\usepackage{tikz}
%\usetikzlibrary{matrix,arrows,calc}
%\tikzset{column sep=3em,row sep=2em,text height=1.5ex, text depth=0.25ex,>=angle 90}
%\tikzset{description/.style={fill=white,inner sep=1pt}}

% hyperliens, 'etiquettes
\usepackage{xr-hyper}
\usepackage{../fichiers-configuration-latex/smfthm}
%\ifx\cacheretiquettes\undefined
%\usepackage[inline]{showlabels} \else \fi % pour rendre visibles les 'etiquettes
%\usepackage{../fichiers-configuration-latex/makeidx} % pour faire un index
%\usepackage[pagebackref,unicode]{hyperref}  % pour voir o`u sont cit'ees les r'ef'erences
\hypersetup{unicode,colorlinks=true,linkcolor=blue,citecolor=Sepia,urlcolor=magenta,
  hypertexnames=false}
%\hypersetup{breaklinks=true} %permet le retour `a la ligne dans les liens trop longs. Incompatible avec pdflatex
%\renewcommand*{\backref}[1]{{\small [cit\'e page(s)~#1.]}}

\usepackage{xspace}
\usepackage{pifont} % pour expose XVII

\renewcommand{\thefootnote}{\roman{footnote}} % pour mettre les notes en bas de page sous la forme (ii) et non (2), qui ressemble trop (?) a un exposant.

\newcounter{pageoriginal}
\newcommand{\pageoriginale}{\refstepcounter{pageoriginal}\marginpar{\footnotesize\textbf{\thepageoriginal}}\xspace}
\let\pageoriginaled\pageoriginale

\usepackage{ifthen}
\newboolean{original}
\def\sisi{\ifthenelse{\boolean{original}}}
\setboolean{original}{false}

\newcounter{cnote}
\def\nde#1{\ifthenelse{\boolean{original}}{}{\stepcounter{cnote}\def\thefootnote{\arabic{cnote}}\protect\footnote[\value{cnote}]{N.D.E.: #1}}}
\def\ndemark{\ifthenelse{\boolean{original}}{}{\stepcounter{cnote}\def\thefootnote{\arabic{cnote}}\protect\footnotemark[\value{cnote}]}}
\def\ndetext#1{\ifthenelse{\boolean{original}}{}{\def\thefootnote{\arabic{cnote}}\protect\footnotetext[\value{cnote}]{N.D.E.: #1}}}

\setmainlanguage{french}
\setotherlanguage{english}

% unicode-math deliberately disables the legacy \Arrowvert glyph.  The
% source uses it once as a rotated equality marker; \Vert has the intended
% appearance and keeps the historical diagram buildable with LuaLaTeX.
\renewcommand{\Arrowvert}{\Vert}


%%
%% Pour afficher le numéro de version Git (exige que LuaTeX soit appelé avec --shell-escape)
%%

\catcode`\%=12
\directlua{
function command_output(cmd)
  local output_pipe = io.popen(cmd)
  local line_iter = output_pipe and output_pipe:lines()
  local line = line_iter and line_iter()
  if output_pipe then output_pipe:close() end
  return line
end
local git_command1 = [[git log --pretty=format:"%h %ai" -1]]
local git_command2 = [[if git diff --quiet HEAD ; then echo clean ; else echo dirty ; fi]]
function git_status()
  local status = command_output(git_command1)
  if status then
    tex.tprint({luatexbase.catcodetables.CatcodeTableLaTeX, string.char(0x5c).."texttt{"}, {-2,status}, {luatexbase.catcodetables.CatcodeTableLaTeX, "}"})
    if not (command_output(git_command2) == "clean") then tex.sprint(-2, " (modifications non validées)") end
  else
    tex.sprint(-2, "?")
  end
end
}
\catcode`\%=14
\def\showgitstatus{\directlua{git_status()}}


% macros pour LaTeX
\renewcommand{\mathsf}{}

% Police différente pour U+2113 SCRIPT SMALL L
\setmathfont[range="2113]{latinmodern-math.otf}

% Give every explicit SGA label its own destination and preserve the printed
% number supplied as the optional argument.  The former definition inherited
% whichever counter anchor happened to precede it and ignored #1 entirely.
\makeatletter
\def\etiquette[#1]#2{%
  \phantomsection
  \protected@edef\@currentlabel{#1}%
  \label{#2}%
}
\makeatother
\def\eqhyperref[#1]#2{(\hyperref[#1]{#2})}
\def\pagehyperref[#1]#2{(page \hyperref[#1]{#2})}
\def\hyperrefext[#1][#2]#3{\hyperref[#2]{#1 #3}}
\def\eqhyperrefext[#1][#2]#3{\hyperref[#2]{#1 #3}}     %cf. exposes 17
\def\neweqhyperrefext[#1][#2]#3{\hyperref[#2]{#1 #3}}  %cf. exposes 17

\def\shat#1{\protect\widehat{#1}}
\def\stilde#1{\protect\widetilde{#1}}
%%%%%%%%%%%%%%%%%%


\def\etcn{etc.\xspace}
\def\etco{etc...\xspace}
\def\etc{\sisi{\etco}{\etcn}}
\def\ldotsn{\sisi{...}{\etcn}}

\def\cf{{cf.}\kern.3em}
\def\Cf{{Cf.}\kern.3em}
\def\ie{{i.e.}\kern.3em}
\def\iev{{i.e.},\xspace}
\def\loccito{loc.\kern2pt cit.\xspace}
\def\loccitn{\textit{loc.\kern2pt cit.}\xspace}
\def\loccit{\sisi{\loccito}{\loccitn}}
\def\EGA{\textup{EGA}\kern.2em}
\def\SGA{\textup{SGA}\kern.3em}

\def\quoi{}

\def\quav{\quad,\quad}
\def\quaet{\quad \text{et}\quad}
\def\fibreproduct#1#2#3{#1 ×_{#3} #2}
\let\fprod\fibreproduct
\def\ccheck#1#2{\mbox{${\check{\check{\rm #1}}}^{\lower2pt\hbox{$\scriptstyle #2$}}$}}
\def\sheaf#1{\protect\mathcal{#1}}
\def\vieuxsheaf#1{{#1}} % cf. problème avec sheaf{a\Phi_i} par exemple...
\DeclareMathOperator{\Hom}{Hom}
\DeclareMathOperator{\can}{can}
\DeclareMathOperator{\Ind}{Ind}
\DeclareMathOperator{\ind}{ind}
%\DeclareMathOperator{\Hhomtop}{\mathcal{H}\!\mathit{omtop}}
\DeclareMathOperator{\Hhomtop}{\mathbf{Homtop}}
\DeclareMathOperator{\Point}{Point}
\DeclareMathOperator{\Pro}{Pro}
\DeclareMathOperator{\Top}{Top}
\DeclareMathOperator{\Mod}{Mod}
\DeclareMathOperator{\Cart}{Cart}
%\DeclareMathOperator{\Mmorsite}{\mathcal{M}\!\mathit{orsite}}
\DeclareMathOperator{\Mmorsite}{\mathbf{Morsite}}
\DeclareMathOperator{\loc}{loc}
\DeclareMathOperator{\Loc}{Loc}
\DeclareMathOperator{\Pic}{Pic}
\DeclareMathOperator{\Et}{Et}
\DeclareMathOperator{\card}{card}
\DeclareMathOperator{\Ens}{Ens}
\DeclareMathOperator{\h}{\underline{H}}
\DeclareMathOperator{\Fonct}{Fonct}
\DeclareMathOperator{\Ker}{Ker}
\DeclareMathOperator{\but}{but}
\DeclareMathOperator{\Ob}{Ob}
\DeclareMathOperator{\ob}{ob}
\DeclareMathOperator{\Lib}{Lib}
\DeclareMathOperator{\esj}{esj}
\DeclareMathOperator{\id}{id}
\DeclareMathOperator{\Id}{Id}
\DeclareMathOperator{\ad}{ad}
\DeclareMathOperator{\Iim}{Im}
\DeclareMathOperator{\pr}{pr}
\DeclareMathOperator{\Cat}{Cat}
\DeclareMathOperator{\car}{car}
\DeclareMathOperator{\SysInd}{SysInd}
\DeclareMathOperator{\Sup}{Sup}
\DeclareMathOperator{\Coker}{Coker}
\DeclareMathOperator{\Filt}{Filt}
\DeclareMathOperator{\acc}{acc}
\DeclareMathOperator{\sur}{sur}
\DeclareMathOperator{\cat}{cat}
\DeclareMathOperator{\Ord}{Ord}
\DeclareMathOperator{\Gr}{Gr}
\DeclareMathOperator{\Cov}{Cov}
\DeclareMathOperator{\Caf}{Caf}
\DeclareMathOperator{\coker}{coker}
\DeclareMathOperator{\coim}{coim}
\DeclareMathOperator{\im}{im}
\DeclareMathOperator{\Ouv}{Ouv}
\DeclareMathOperator{\diag}{diag}
\DeclareMathOperator{\Esp}{Esp}
\DeclareMathOperator{\TOP}{TOP}
\DeclareMathOperator{\Unif}{Unif}
\DeclareMathOperator{\Iind}{\mathbf{Ind}}
\DeclareMathOperator{\Hhomcart}{\mathbf{Homcart}}
\DeclareMathOperator{\Hhomtopess}{\mathbf{Homtopess}}
\DeclareMathOperator{\Ppoints}{\mathbf{Points}}
\DeclareMathOperator{\Ppoint}{\mathbf{Point}}
\DeclareMathOperator{\Ttors}{\mathbf{Tors}}
\DeclareMathOperator{\Ssex}{\mathbf{Sex}}
\DeclareMathOperator{\Ffib}{\mathbf{Fib}}
\DeclareMathOperator{\Vvois}{\mathbf{Vois}}
\DeclareMathOperator{\Ppointess}{\mathbf{Pointess}}
\DeclareMathOperator{\scrHom}{\mathbf{Hom}}
\DeclareMathOperator{\sheafPic}{\mathbf{Pic}}
\DeclareMathOperator{\sheafAut}{\mathbf{Aut}}
\DeclareMathOperator{\Hhomtoplocan}{\mathbf{Homtoplocan}}
\DeclareMathOperator{\Hhomtopan}{\mathbf{Homtopan}}
\DeclareMathOperator{\Pptgeom}{\mathbf{Ptgéom}}
\DeclareMathOperator{\Morsite}{Morsite}
\DeclareMathOperator{\Fib}{Fib}
\DeclareMathOperator{\Fais}{Fais}
\DeclareMathOperator{\kar}{kar}
\DeclareMathOperator{\Groupes}{Groupes}
\DeclareMathOperator{\Restr}{Restr}
\DeclareMathOperator{\Isom}{Isom}
\DeclareMathOperator{\Prol}{Prol}
\DeclareMathOperator{\Res}{Res}
\DeclareMathOperator{\Tors}{Tors}
\DeclareMathOperator{\esp}{esp}
\DeclareMathOperator{\es}{es}
\DeclareMathOperator{\od}{od}
\DeclareMathOperator{\gr}{gr}
\DeclareMathOperator{\sob}{sob}
\DeclareMathOperator{\dfn}{dfn}
\DeclareMathOperator{\Sch}{Sch}
\DeclareMathOperator{\ouv}{ouv}
\DeclareMathOperator{\cons}{cons}
\DeclareMathOperator{\Nerf}{Nerf}
\DeclareMathOperator{\Points}{Points}
\DeclareMathOperator{\Espsob}{Espsob}
\DeclareMathOperator{\Homtop}{Homtop}
\DeclareMathOperator{\Pt}{Pt}
\DeclareMathOperator{\Quot}{Quot}
\DeclareMathOperator{\glob}{glob}
\DeclareMathOperator{\Homfib}{Homfib}
\DeclareMathOperator{\Qouv}{Qouv}
\DeclareMathOperator{\Quotouv}{Quotouv}
\DeclareMathOperator{\ens}{ens}
\DeclareMathOperator{\ab}{ab}
\DeclareMathOperator{\an}{an}
\DeclareMathOperator{\fs}{fs}
\DeclareMathOperator{\red}{red}
\DeclareMathOperator{\Ab}{Ab}
\DeclareMathOperator{\cosq}{cosq}
\DeclareMathOperator{\Ann}{Ann}
\DeclareMathOperator{\Zar}{Zar}


%%%%%%%%%%%%
%%%%% SGA4-2
\DeclareMathOperator{\Sex}{Sex}
\DeclareMathOperator{\vois}{vois}
\DeclareMathOperator{\Vois}{Vois}
\DeclareMathOperator{\Ext}{Ext}
\DeclareMathOperator{\ext}{ext}
\DeclareMathOperator{\Fer}{Fer}
\DeclareMathOperator{\Add}{Add}
\DeclareMathOperator{\Cocart}{Cocart}
\DeclareMathOperator{\cosk}{cosk}
\DeclareMathOperator{\plat}{plat}

\DeclareMathOperator{\faisc}{faisc}
\DeclareMathOperator{\sk}{sk}
\DeclareMathOperator{\ch}{ch}

\def\HH{\mathbf{H}}
\def\qc{\mathrm{qc}}
\def\et{\mathrm{et}}
\def\coh{\mathrm{coh}}

\def\Hvr#1{%
\begin{array}[b]{l}
{\kern2pt\scriptstyle r}\\[-8pt]
{\kern1pt\scriptstyle \vee}\\[-5pt]
H^{#1}
\end{array}}

\def\Hvv#1{%
\begin{array}[b]{l}
\kern2pt\rotatebox{90}{${\scriptscriptstyle\ll}$}\\[-6pt]
H^{#1}
\end{array}}

\def\Hvi#1{%
\begin{array}[b]{l}
{\kern0.2pt\scriptscriptstyle \infty}\\[-8pt]
{\kern1pt\scriptstyle \vee}\\[-5pt]
H^{#1}
\end{array}}

\def\Hv#1{%
\begin{array}[b]{l}
{\kern1pt\scriptstyle \vee}\\[-5pt]
H^{#1}
\end{array}}

\DeclareMathOperator{\cart}{cart}
\DeclareMathOperator{\Ffaisc}{\mathbf{Faisc}}


\def\Ev#1{%
\begin{array}[b]{l}
{\kern1pt\scriptstyle \vee}\\[-5pt]
E'_{#1}
\end{array}}

\DeclareMathOperator{\Cont}{Cont}
\DeclareMathOperator{\Ccont}{\mathscr{C}\!\mathit{ont}}
\DeclareMathOperator{\Site}{Site}
\newcommand\MM{\mathbf{M}}
\def\limtop#1{{\displaystyle{\varprojlim_{#1}\!\mbox{\rm top }}}}
\def\Limtop{\mathop{\varprojlim\!\mbox{\rm top}}}
\DeclareMathOperator{\spec}{spec}
\def\ss{\scriptstyle{\mathscr{S}}}
\DeclareMathOperator{\Gal}{Gal}
\def\smallm{\scriptstyle{\mathscr{M}}}
\DeclareMathOperator{\End}{End}
\DeclareMathOperator{\ass}{ass}
\DeclareMathOperator{\zar}{zar}
\DeclareMathOperator{\Ogg}{Ogg}
\DeclareMathOperator{\res}{res}
\DeclareMathOperator{\tr}{tr}
\DeclareMathOperator{\Supp}{Supp}
\DeclareMathOperator{\codim}{codim}
\DeclareMathOperator{\rk}{rk}
\DeclareMathOperator{\cur}{cur}
\DeclareMathOperator{\cdqc}{cdqc}
\newcommand\DD{\mathbf{D}}
\DeclareMathOperator{\cl}{cl}
\def\JJ{\mathbf{J}}
\def\FF{\mathbf{F}}
\def\BB{\mathbf{B}}
\DeclareMathOperator{\Alget}{Alget}
\DeclareMathOperator{\Ass}{Ass}
\DeclareMathOperator{\supp}{supp}
\DeclareMathOperator{\rad}{rad}
\DeclareMathOperator{\prof}{prof}
\DeclareMathOperator{\cint}{int}
\DeclareMathOperator{\sic}{sic}
\DeclareMathOperator{\Tor}{Tor}
\DeclareMathOperator{\tors}{tors}
\DeclareMathOperator{\Sym}{Sym}
\DeclareMathOperator{\parf}{parf}
\DeclareMathOperator{\tor}{tor}
\DeclareMathOperator{\defn}{defn}
\DeclareMathOperator{\Div}{Div}
\DeclareMathOperator{\Pol}{Pol}
\DeclareMathOperator{\Fex}{Fex}
\DeclareMathOperator{\Lie}{Lie}
\DeclareMathOperator{\tot}{tot}
\DeclareMathOperator{\Aut}{Aut}
\DeclareMathOperator{\HOM}{HOM}
\DeclareMathOperator{\pch}{pch}
\DeclareMathOperator{\Ch}{Ch}
\DeclareMathOperator{\EXT}{EXT}
\DeclareMathOperator{\PIC}{PIC}
\DeclareMathOperator{\Rf}{Rf}
\DeclareMathOperator{\inj}{inj}
\newcommand\VV{\mathbf{V}}

\def\Xv{%
\begin{array}[b]{l}
{\kern1pt\scriptstyle \vee}\\[-5pt]
X
\end{array}}

%%%%%%%
\newcommand\TT{\mathbf{T}}
\newcommand\XX{\mathbf{X}}
\newcommand\EE{\mathbf{E}}
%%%%%%

\def\qedsymbol{\textup{C.Q.F.D.}}
\let\CQFD\qed
\let\cqfd\qed

%%%%%\bigdot
\font\bigsymb=cmsy10 at 4pt
\def\bigdot{{\kern1.2pt\raise 1.5pt\hbox{\bigsymb\char15}}}
%%%%%

\def\Lright{\mathop{\vcenter{\hbox{\hbox{$\longrightarrow$}\lower8pt \hbox{\kern-16pt\hbox{$\longrightarrow$}}}}}}

\def\Aa{{\mathcal A}}
\def\Bb{{\mathcal B}}
\def\Cc{{\mathcal C}}
\def\Dd{{\mathcal D}}
\def\Ee{{\mathcal E}}
\def\Ff{{\mathcal F}}
\def\Gg{{\mathcal G}}
\def\Hh{{\mathcal H}}
\def\Ii{{\mathcal I}}
\def\Jj{{\mathcal J}}
\def\Kk{{\mathcal K}}
\def\Ll{{\mathcal L}}
\def\Mm{{\mathcal M}}
\def\Nn{{\mathcal N}}
\def\Oo{{\mathcal O}}
\def\Pp{{\mathcal P}}
\def\Qq{{\mathcal Q}}
\def\Rr{{\mathcal R}}
\def\Ss{{\mathcal S}}
\def\Tt{{\mathcal T}}
\def\Uu{{\mathcal U}}
\def\Vv{{\mathcal V}}
\def\Ww{{\mathcal W}}
\def\Xx{{\mathcal X}}
\def\Yy{{\mathcal Y}}
\def\Zz{{\mathcal Z}}

\newcommand{\CC}{\mathbf{C}}
\newcommand{\GG}{\mathbf{G}}
\newcommand{\KK}{\mathbf{K}}
\newcommand{\LL}{\mathbf{L}}
\newcommand{\NN}{\mathbf{N}}
\newcommand{\PP}{\mathbf{P}}
\newcommand\QQ{\mathbf{Q}}
\newcommand\RR{\mathbf{R}}
\newcommand\ZZ{\mathbf{Z}}

\def\to{\mathchoice{\longrightarrow}{\rightarrow}{\rightarrow}{\rightarrow}}
\def\from{\mathchoice{\longleftarrow}{\leftarrow}{\leftarrow}{\leftarrow}}
\def\mto{\mathchoice{\longmapsto}{\mapsto}{\mapsto}{\mapsto}}
\def\mfrom{\mathrel{\from\joinrel\mapstochar\,}}
\def\hto{\mathrel{\lhook\joinrel\to}}
\def\hfrom{\mathrel{\from\joinrel\rhook}}

\def\implies{\mathchoice{\Longrightarrow}{\Rightarrow}{\Rightarrow}{\Rightarrow}}
\def\implied{\mathchoice{\Longleftarrow}{\Leftarrow}{\Leftarrow}{\Leftarrow}}
\def\iff{\mathchoice{\Longleftrightarrow}{\Leftrightarrow}{\Leftrightarrow}{\Leftrightarrow}}

\def\To#1{\mathchoice{\xrightarrow{\labelstyle~~#1~}}{\stackrel{#1}{\longrightarrow}}{}{}}
\def\Mto#1{\mathrel{\mapstochar\To{#1}}}
\def\From#1{\mathchoice{\xleftarrow{\labelstyle~~#1~}}{\stackrel{#1}{\longleftarrow}}{}{}}

\newcommand{\isomto}{\overset{\sim}{\longrightarrow}}
\newcommand{\isomfrom}{\overset{\sim}{\longleftarrow}}

\newcommand{\IFF}{\unskip~$\iff$~\ignorespaces}
\newcommand{\THEN}{\unskip~$\implies$~\ignorespaces}

\let\leq\leqslant
\let\geq\geqslant
\let\le\leqslant
\let\ge\geqslant

\def\cf{{cf.}\kern.3em}
\def\Cf{{Cf.}\kern.3em}
\def\ie{{i.e.}\kern.3em}
\def\iev{{i.e.},\xspace}
%\def\loccito{loc.\kern2pt cit.\xspace}
%\def\loccitn{\textit{loc.\kern2pt cit.}\xspace}
%\def\loccit{\sisi{\loccito}{\loccitn}}
\def\EGA{\textup{EGA}\kern.2em}
\def\SGA{\textup{SGA}\kern.3em}

\let\oldS\S
\def\S{\oldS\kern .15em}
\def\ptbl{.\kern .2em}
\def\resp{resp.\kern.3em}
\def\Exp{Exp.\kern.3em}

\let\bar\overline
\let\hat\widehat
\let\tilde\widetilde
\let\dpl\displaystyle
\let\ldots\dots
\let\emptyset\varnothing
\let\skpt\skippointrait
\let\epsilon\varepsilon

\def\sfrac#1#2{{#1}/{#2}}

\DeclareMathOperator{\cd}{cd}
%\DeclareMathOperator{\Hhom}{\mathcal{H}\!\mathit{om}}
\DeclareMathOperator{\Hhom}{\mathbf{Hom}}
%\DeclareMathOperator{\Eext}{\mathcal{E}\!\mathit{xt}}
\DeclareMathOperator{\Eext}{\mathbf{Ext}}
\DeclareMathOperator{\Spec}{Spec}
\DeclareMathOperator{\Tr}{Tr}

\let\emptyset=\varnothing

%%%%%%%%%%%%%%%
%%% expose XVII
%%%%%%%%%%%%%%%

\makeatletter
\let\old@optexpose\optexpose
\newcommand{\empty@optexpose}[1]{}
\newcommand{\paren@optexpose}[1]{\old@optexpose{#1}(}
\newcommand{\refnoexpose}[1]{\let\optexpose\empty@optexpose\ref{#1}\let\optexpose\old@optexpose}%\refnoexpose{X.cor5.2} engendre "5.2" au lieu de "X 5.2"
\newcommand{\neweqref}[1]{\let\optexpose\paren@optexpose\textup{\ref{#1})}\let\optexpose\old@optexpose}%\neweqref{XII.eq4.1} engendre "XII (4.1)" au lieu de "(XII 4.1)"
\makeatother
\newcommand{\refvb}[2][VI${}_{\text{B}}$ ]{\sisi{#1\refnoexpose{#2}}{\ref{#2}}}

\newcommand{\leftexp}[2]{{\vphantom{#2}}^{#1}{#2}}
\newcommand{\AAA}{\mathbf{A}}
\newcommand{\Sf}{\mathfrak{S}}
\DeclareMathOperator*{\quotprojlim}{\text{``$\varprojlim$''}}
\DeclareMathOperator*{\quotinjlim}{\text{``$\varinjlim$''}}
\DeclareMathOperator{\Fl}{F\sisi{\ell}{l}}
\DeclareMathOperator{\Pl}{P\sisi{\ell}{l}}
%\DeclareMathOperator{\cl}{c\sisi{\ell}{l}}
\DeclareMathOperator{\tdf}{tdf}
\DeclareMathOperator{\rg}{rg}
\newcommand{\iInd}{\sisi{I}{i}nd}
\newcommand{\iind}{ind}
\newcommand{\ppro}{pro}
\newcommand{\tTor}{\sisi{T}{t}or}
\newcommand{\ttor}{tor}
\newcommand{\chap}{Chap.}
\newcommand{\Ffex}{\Ff\mathit{ex}}
\newcommand{\sheafk}{\underline{k}}
\newcommand{\sheafn}{\underline{n}}
\newcommand{\sheafGamma}{\underline{\Gamma}}

\sisi{
\DeclareMathOperator{\simarrow}{\rule{.15mm}{1.8mm}\kern -3.4pt\rightsquigarrow}
\newcommand{\sheafell}{\underline{\ell}}
\newcommand{\sheafZZ}{\underline{\ZZ}}
}{
\newcommand{\nderef}[1]{la N.D.E.~(\ref{#1}) page~\pageref{#1}}
}

%%%%%%%%%%%%%%%%%
%%%%% expose XVII
%%%%%%%%%%%%%%%%%



\externaldocument{../01/01.tex}
\externaldocument{../02/02.tex}
\externaldocument{../03/03.tex}
\externaldocument{../04/04.tex}
\externaldocument{../05/05.tex}
\externaldocument{../05bis/05bis.tex}
\externaldocument{../06/06.tex}
\externaldocument{../07/07.tex}
\externaldocument{../08/08.tex}
\externaldocument{../09/09.tex}
\externaldocument{../10/10.tex}
\externaldocument{../11/11.tex}
\externaldocument{../12/12.tex}
%\externaldocument{../13/13.tex}
\externaldocument{../14/14.tex}
\externaldocument{../15/15.tex}
\externaldocument{../16/16.tex}
\externaldocument{../17/17.tex}
\externaldocument{../18/18.tex}
\externaldocument{../19/19.tex}

\setboolean{original}{false}

\begin{document}

\begin{center}
XIII. Théorème de changement de base pour un morphisme propre : fin de la démonstration\\
M. Artin 
\end{center}

\selectlanguage{english}version : \showgitstatus\selectlanguage{french}
\tableofcontents

\else
\chapter{Théorème de changement de base pour un morphisme propre : fin de la démonstration}
%\addtocontents{toc}{M. Artin}
\begin{center}
M. Artin 
\end{center}
\fi

\setcounter{pageoriginal}{131}

\section{Le cas projectif et plat}\etiquette[1]{XIII.sec1}\pageoriginale

Rappelons l'énoncé de \hyperrefext[XII][XII.thm5.9]{5.9}. bis dans le cas envisagé ici:

\begin{enonce*}{Proposition 1.1}\etiquette[1.1]{XIII.prop1.1}
Soit $S$ le spectre d'un anneau hensélien noethérien, et
$f:X\rightarrow S$ un morphisme projectif et plat. Soit $X_{\circ}$ la
fibre fermée de $X/S$. Le foncteur de restriction
$$
\Et(X)\longrightarrow \Et(X_{\circ})
$$
est une équivalence de catégories.
\end{enonce*}

Comme nous l'avons déjà remarqué, il résulte de \hyperrefext[XII][XII.prop6.5]{6.5} (i) que la
flèche est pleinement fidèle. Il reste à démontrer que chaque
revêtement étale $Y_{\circ}$ de $X_{\circ}$ est induit par un
revêtement étale $Y$ de $X$.

\begin{enonce*}{Lemme 1.2}\etiquette[1.2]{XIII.lem1.2}
Soit $S$ localement noethérien, $f:X\rightarrow S$ projectif et plat,
et $Y\subset X$ un sous-schéma. Le sous-foncteur
$\epsilon=\prod_{X/S}Y/X$: $(\Sch)/S(\Ens)$ du foncteur final « qui exprime la condition $Y=X$ », \ie tel que pour $S'\rightarrow
S$, on ait
$$
\epsilon(S')=
\begin{cases}
∅ & \mbox{si }\quad Y'\neq X'\\
\{∅\} &\mbox{si }\quad Y'=X'
\end{cases}
$$
(où\pageoriginale l'on désigne par un prime l'effet du changement de base
$S'\rightarrow S$) est représentable par un sous-préschéma $Z$ de
$S$. Si $Y$ est ouvert (\resp fermé), il en est de même de $Z$.
\end{enonce*}

\noindent
\emph{Démonstration}. Soit $Y$ fermé dans l'ouvert $U$ de $X$, et soit
$C$ le fermé complémentaire de $U$. Puisque $F$ est propre, l'image de
$C$ dans $S$ est un fermé $D$, et il est évident que la condition
$U=X$ est représentée par l'ouvert $S-D$ de $S$. On peut donc supposer
qu'on a déjà $U=X$, \ie que $Y$ est un sous-schéma fermé. Considérons
le morphisme surjectif $\Oo_{X}\xrightarrow{\alpha}\Oo_{Y}$ de
$\Oo_X$-modules; dire que $Y=X$ revient au même que de dire qu'il
existe un morphisme $\Oo_Y\xrightarrow{\beta}\Oo_X$ tel que le
composé $\beta\alpha$ soit l'identité, c'est-à-dire, qu'il existe une
section de $\Hhom(\Oo_Y,\Oo_X)$ qui relève la section identique
de $\Hhom(\Oo_X,\Oo_X)$. Or il résulte de \EGA III 7.7.8 que le
foncteur qui à $S'/S$ associe l'ensemble des sections de
$\Hhom(\Oo_Y,\Oo_X)$ (\resp $\Hhom(\Oo_X,\Oo_X)$) est
représenté par un fibré vectoriel $V_1$ (\resp $V_2$) sur $S$. Soient
$\varphi:V_1\rightarrow V_2$ le morphisme $u\mapsto u\alpha$ induit
par $\Oo_X\rightarrow \Oo_Y$, et $s:S\rightarrow V_2$ la section
correspondant à l'identité $\Oo_X\rightarrow \Oo_X$. L'image
inverse de la section $s$ par $\varphi$ s'identifie à un sous-schéma
fermé $Z$ de $S$, et il est clair que $Z$ représente le foncteur
$\epsilon$, d'où le lemme.

\begin{enonce*}{Lemme 1.3}\etiquette[1.3]{XIII.lem1.3}
Soient $S$ un schéma localement noethérien, $f:X\rightarrow S$ un
morphisme projectif et plat, et $E$ un faisceau localement libre sur
$X$. Soit
$$
QE:(\Sch)/S\longrightarrow (\Ens)
$$
le foncteur qui à $S'/S$ associe l'ensemble $QE(S')$ des faisceaux
quotients $A'$ de $E'=E\otimes_{O_S}\Oo_{S'}$ munis d'une structure
d'Algèbre étale sur $\Oo_{X'}$. Ce foncteur est représentable par un
schéma localement de type fini sur $S$.
\end{enonce*}

\noindent
\emph{Démonstration}.\pageoriginale Le foncteur $\sheaf{\Quot}(E)$,
$\sheaf{\Quot}(E)(S')=$ ensemble des quotients de $E'$ plat au-dessus de
$S'$, est représentable par un schéma localement de type fini (TDTE
IV, exp.~221), disons par $Z/S$. Soit $M_u$ le quotient
« universel » de $E\otimes_{\Oo_S}\Oo_Z$ au-dessus de $X_Z$. Le
sous-foncteur de $\sheaf{\Quot}(E)$ des quotients que sont plats
\emph{au-dessus de} $X'$ est représenté par un ouvert $U$ de $Z$. En
effet, soit $M'$ un quotient de $E'$, donné par un $S$-morphisme
$S'\rightarrow Z$. Puisque $M'$ est plat au-dessus de $S'$, dire qu'il
est plat au-dessus de $X'$ revient au même que de dire que pour chaque
point $s'\in S'$, le module $M'_{s'}$ est plat au-dessus de la fibre
$X'_{s'}$ (\EGA IV 5.9), ce qui équivaut à $(M_u)_z$ plat au-dessus de
la fibre $Z_z$ au point $z$ image de $s'$. Or l'ensemble $U$ des $z\in
Z$ tels que $(M_u)_z$ soit plat sur $Z_z$, \ie que $M_u$ soit plat sur
$Z$ en les points de $Z_s$, est ouvert dans $Z$ en vertu de \EGA IV
11.1.5, d'où l'assertion.

Il suffit maintenant de représenter le foncteur relatif $QE/U[f]$. On
est donc réduit au lemme suivant, en remplaçant $S$ par $U$ et $X$
par $\fprod{X}{U}{S}$:

\begin{enonce*}{Lemme 1.4}\etiquette[1.4]{XIII.lem1.4}
Soient $S$ un schéma localement noethérien, $f:X\rightarrow S$ un
morphisme projectif et plat, et $M$ un faisceau localement libre sur
$\Oo_X$. Le foncteur $\Alget:(\Sch)/S\rightarrow (\Ens)$, qui à
$S'/S$ associe l'ensemble $\Alget(S')$ des structures d'algèbre sur
$M'$ étales au-dessus de $\Oo_{X'}$, est représentable par un schéma
de type fini sur $S$.
\end{enonce*}

\noindent
\emph{Démonstration}. Nous dirons, pour abréger,
« représentable » en sous-entendant « par un schéma relatif
de \emph{type fini} ». Les fibrés vectoriels rencontrés seront tous définis par des
faisceaux \emph{cohérents}, donc seront de type fini. Une loi de
composition bilinéaire dans $M'$ est donnée par une section du
faisceau localement libre $\Hhom$ $(M'\otimes M', M')$. Puisque
$\Hhom$ $(M'\otimes M', M')$ est $S$-plat, il résulte
de\pageoriginale \EGA III 7.7.6 que le foncteur « loi de composition
bilinéaire » est représentée par un fibré vectoriel $Z$ sur $S$. Il
faut vérifier que les conditions supplémentaires sur la loi d'être
associative, commutative, et d'avoir une section unité, sont
représentées par un sous-schéma de $Z$, et on peut (en considérant le
foncteur relatif) remplacer $S$ par $Z$, donc peut supposer qu'on a
une loi de composition donnée sur $M$, et montrer qu'on peut
représenter le sous-foncteur de $S$ exprimant les conditions
précédentes par un sous-schéma de $S$.

Or pour l'associativité, la condition est que deux sections de
$\Hhom$ $(M\otimes M\otimes M, M)$ deviennent égales par changement
de base $S'\rightarrow S$. Soit $Y$ le produit fibré de ces sections
au-dessus du fibré vectoriel correspondant à $\Hhom$ $(M\times
M\otimes M, M)$. Le schéma $Y$ s'identifie à un sous-schéma fermé de
$X$, et l'associativité après $S'\rightarrow S$ est exprimé par la
condition $Y'=X'$. Elle définit donc un sous-foncteur de $S$
représentable par un sous-schéma fermé de $S$ d'après \hyperref[XIII.lem1.2]{1.2}. Pour la
commutativité la condition est que deux sections de $\Hhom$
$(M\otimes M,M)$ deviennent égales, ce qui définit encore un
sous-foncteur de $S$ représentable par un sous-schéma fermé d'après
1.2. Supposons enfin, ce qui est maintenant loisible, que la loi est
déjà associative et commutative, et examinons la condition pour qu'il
existe une section unité, qui sera alors unique. Le foncteur qui en
$S'/S$ a comme valeur l'ensemble des sections de $M'$ est
représentable par un fibré vectoriel, d'après \EGA III 7.7.6, et tout
revient à représenter le sous-foncteur de ce dernier correspondant aux
\emph{sections unités} (compte tenu de l'unicité qu'on vient de
rappeler). Relativisant comme d'habitude, on peut supposer qu'une
section $\epsilon$ de $M$ est déjà donnée, et montrer que le
sous-foncteur de $S$ exprimant qu'elle devient une\pageoriginale
section unité est représentable. Or, la condition que $\epsilon$ soit
une unité est que le composé des morphismes évidents $M'\rightarrow
M'\otimes M'\rightarrow M'$ soit l'identité, c'est-à-dire que deux
sections déterminées de $\Hom$ $(M,M)$ deviennent égales. On conclut
encore grâce à~\hyperref[XIII.lem1.2]{1.2}.

Supposons enfin qu'une structure d'algèbre commutative unitaire soit
donnée sur le module localement libre $M$. La condition d'être étale
est ouverte dans $X$ (\SGA I 4.5), donc représentable sur $S$ par un
sous-schéma (ouvert). Ceci achève la démonstration de \hyperref[XIII.lem1.4]{1.4}, donc de
\hyperref[XIII.lem1.3]{1.3}.

\setcounter{subsection}{4}
\subsection{}\etiquette[1.5]{XIII.subsec1.5}
Soit maintenant $Z/S$ l'objet qui représente le foncteur envisagé dans
\hyperref[XIII.lem1.3]{1.3}, et examinons la condition sur $Z$ d'être lisse sur $S$ en un
point $z$. En traduisant \SGA 1 III 3.1 en termes du foncteur $QE$, on
trouve la condition suivante:

Soient $B$ un anneau artinien local, $\Mm=rB$, $J$ un idéal de
$B$ tel que $J\Mm=0$, $S'=\Spec B$, $S''=\Spec B/J$ et
$S'\rightarrow S$ un morphisme. On utilise des notations évidentes
pour les effets des changements de bases $S'\rightarrow S$ et
$S''\rightarrow S$.

\begin{itemize}
\item[$(\ast)$] Pour chaque quotient $A''$ de $E''$ muni d'une
structure d'algèbre étale sur $\Oo_{X''}$, correspondant à un
$S$-morphisme $S''\rightarrow Z$ d'image $z$, il existe un quotient
$A'$ de $E'$, muni d'une structure d'algèbre étale sur $\Oo_{X'}$,
qui l'induise.
\end{itemize}

Or puisque $X'$ et $X''$ ont même espace sous-jacent, une algèbre $A'$
qui induit $A''$ est déjà déterminée à isomorphisme unique près (\SGA
I 8.3). On\pageoriginale peut donc exprimer la condition
$({\ast}{\ast})$ ainsi:

\begin{itemize}
\item[$({\ast}{\ast})$] Soit $A'$ une algèbre étale sur $\Oo_{X'}$,
et $A''$ l'algèbre induite de $A'$ sur $X''$. Chaque homomorphisme
surjectif
$$
E''\xrightarrow{\varphi''}A''
$$
correspondant à un morphisme $S''\rightarrow Z$ d'image $z$, se relève
en un homomorphisme (nécessairement surjectif en vertu du lemme de
Nakayama)
$$
E'\xrightarrow{\varphi'}A'{\quoi}.
$$
\end{itemize}

De la suite exacte
$$
0\longrightarrow J\longrightarrow B\longrightarrow B/J\longrightarrow
0
$$
on déduit une suite exacte de faisceaux sur $X'$:
$$
0\longrightarrow J\otimes \Hhom(E',A')\longrightarrow
\Hhom(E',A')\longrightarrow \Hhom(E'', A'')\longrightarrow
0{\quoi}.
$$
L'obstruction au relèvement de $\varphi''$ se trouve dans $H^1$
$(X',J\otimes \Hhom(E',A'))$. Soit $S_{\circ}=\Spec
B/\Mm$. Puisque $J\Mm=0$, on a
$$
J\otimes \Hhom(E',A')\simeq J\otimes
\Hhom(E_{\circ},A_{\circ}){\quoi},
$$
et il est clair que
$$
H^1(X',J\otimes \Hhom(E_{\circ},A_{\circ}))\simeq J\otimes
H^1(X_{\circ},\Hhom(E_{\circ},A_{\circ})){\quoi}.
$$
Tenant compte du fait que la cohomologie commute à l'extension des
corps de base, on trouve le

\begin{enonce*}{Corollaire 1.6}\etiquette[1.6]{XIII.cor1.6}
Soient $Z/S$ le schéma qui représente le foncteur de $1.3$,
$X_Z=\fprod{X}{Z}{S}$, et $A_Z$ l'Algèbre étale quotient
(« générale ») de $E_Z$. Alors $Z/S$ est lisse en tout point $z_{\circ}$ tel que
$$
H^1(X_{\circ},\Hhom(E_{\circ},A_{\circ}))=0{\quoi},
$$
où\pageoriginale l'indice zéro désigne la restriction à la fibre
$X_{\circ}$ de $X_Z/Z$ au point $z_{\circ}$.
\end{enonce*}

Nous pouvons maintenant achever la démonstration du théorème \hyperref[XIII.prop1.1]{1.1}: avec
les notations du théorème, soit $Y_{\circ}/X_{\circ}$ un revêtement
étale, donné par un faisceau d'Algèbres $A_{\circ}$, étale sur
$\Oo_{X_{\circ}}$. Pour $n$ assez grand, $A_{\circ}(n)$ est engendré
par ses sections (\EGA III 2.2.1). On a donc une surjection
$$
\Oo^N_{X_{\circ}}\longrightarrow A_{\circ}(n)\longrightarrow
0{\quoi},
$$
d'où une surjection
$$
O^N_{X_{\circ}}(-n)=E_{\circ}\longrightarrow A_{\circ}\longrightarrow
0{\quoi}.
$$
Choisissons $n$ assez grand pour qu'on ait aussi
\begin{equation*}
H^1(X_{\circ},\Hhom(E_{\circ}, A_{\circ}))=0{\quoi},\tag{$+$}
\end{equation*}
ce qui est possible en vertu de (\EGA III 2.2.1), puisque
$$
\Hhom(E_{\circ},A_{\circ})=[\Hhom(\Oo^N_{X_{\circ}},A_{\circ})](n)=
[A_{\circ}(n)]^N{\quoi}.
$$
Posons $E=O^N_X(-n)$, et soit $Z/S$ le schéma qui représente le
foncteur envisagé dans \hyperref[XIII.lem1.3]{1.3}. Le quotient $A_{\circ}$ de $E_{\circ}$
correspond à un point $z_{\circ}\in Z$ de la fibre fermée de $Z/S$,
rationnel sur $k(s)$, et d'après \hyperref[XIII.subsec1.5]{1.5} et $(+)$, $Z/S$ est lisse au
point $z_{\circ}$. Or puisque $S$ est le spectre d'un anneau
\emph{hensélien} et $Z/S$ est lisse au point
$z_{\circ}$ de la fibre fermée, on peut trouver une section de $Z/S$ passant par $z_{\circ}$: c'est par exemple
immédiat à partir de la définition originale (\SGA II 1.1) de la
lissité. Il s'ensuit que $A_{\circ}$ est induit par une Algèbre $A$
étale sur $X$ (quotient de $E$), ce qui achève la démonstration de
\hyperref[XIII.prop1.1]{1.1}.

\section{Le cas de dimension relative $≤ 1$}\etiquette[2]{XIII.sec2}\pageoriginale 

L'assertion est la suivante:

\begin{enonce*}{Proposition 2.1}\etiquette[2.1]{XIII.prop2.1}
Avec les notations de \rm{\hyperrefext[XII][XII.thm5.9]{5.9}} bis, le foncteur
$\Et(X)\rightarrow \Et(X_{\circ})$ est une équivalence de catégories
si $S$ est noethérien strictement local, et $f$ est projectif et de
dimension $\leq 1$.

D'après {\rm \hyperrefext[XII][XII.lem8.5]{8.5}}, cela achèvera la démonstration de {\rm
  \hyperrefext[XII][XII.thm5.1]{5.1} 
(ii)}, qui généralise à la fois \hyperref[XIII.prop1.1]{1.1} et \hyperref[XIII.prop2.1]{2.1}.
\end{enonce*}

\begin{enonce*}{Lemme 2.2}\etiquette[2.2]{XIII.lem2.2}
Soit $S$ le spectre d'un anneau local noethérien $A$, et soit $X/S$
projectif, de dimension relative $\leq n$. Il existe un $S$-schéma
projectif et plat $Z/S$ de dimension relative $\leq n$ et une
immersion fermée $X\hookrightarrow Z$ au-dessus de $S$.
\end{enonce*}

\noindent
\emph{Démonstration}. Soit $X\rightarrow \PP^N_S$ un plongement
projectif au-dessus de $S$, et soit $J\subset A[x_{\circ},\ldots,x_N]$
l'idéal homogène qui définit $X$. Soient $k=A/rA$, $S_{\circ}=\Spec
k$, $X_{\circ}=X\otimes_S k$. Alors $X_{\circ}$ est défini par l'idéal
homogène $J_{\circ}=\Iim(J\rightarrow k[X_{\circ},\ldots,X_N])$. Comme
$\dim X_{\circ}\leq n$, il est bien connu qu'il existe des éléments
$f^{\circ}_1,\ldots,f^{\circ}_{N-n}$ de $J_{\circ}$ qui définissent un
sous-schéma $Z_{\circ}$ de $\PP^N_k$ de dimension $n$. Soient
$f_i$ $(0 \leq i\leq N-n)$ des éléments de $J$ qui relèvent les
$f^{\circ}_i$, et soit $Z$ le sous-schéma de $\PP^N_{\bar{S}}$
défini par $f_1,\ldots,f_{\nu}$. En vertu de \EGA IV 11.3, $Z$ est plat
au-dessus de $S$, d'autre part en vertu de \EGA IV 13.1.5 il est de
dimension relative $\dim Z_{\circ}=n$. Cela achève de prouver \hyperref[XIII.lem2.2]{2.2}.

\noindent
\emph{Démonstration de} \hyperref[XIII.prop2.1]{2.1}. Soit $X/S$ projectif de dimension
relative $\leq 1$, avec $S$ noethérien strictement local. Nous pouvons
supposer $X$ (donc $X_{\circ}$) connexe. Soit $X\rightarrow Z$ une
immersion fermée au-dessus de $S$ avec $Z$ projectif, plat et de
dimension $\leq 1$, \cf \hyperref[XIII.lem2.2]{2.2}. Soit $s_{\circ}$ le point fermé de
$S$. On a un diagramme correspondant\pageoriginale de schémas
$$
\xymatrix@C=2.5cm@R=1.25cm{
\Et(Z)  \ar[r] \ar[d]^{a}& \Et(X) \ar[d]^{b}                  \\
\Et(Z_0)  \ar[r]^{c} & \Et(X_{0}).
}
$$
On veut démontrer que $b$ est une équivalence. D'après \hyperrefext[XII][XII.lem5.8]{5.8}, $b$
est pleinement fidèle, et il reste à démontrer qu'il est
essentiellement surjectif. Or $a$ est essentiellement surjectif d'après
\hyperref[XIII.prop1.1]{1.1}, et il suffit donc de démontrer que $c$ est essentiellement
surjectif. On est ainsi réduit (laissant tomber le ${}_{\circ}$) à la
proposition suivante:

\begin{enonce*}{Proposition 2.3}\etiquette[2.3]{XIII.prop2.3}
Soient $k$ un corps séparablement clos, $Z$ un $k$-schéma, localement
de type fini et de dimension $\leq 1$, et $X\hookrightarrow Z$ un
sous-schéma fermé connexe. Alors tout revêtement étale de $X$ est
induit par un revêtement étale de $Z$.
\end{enonce*}

\noindent
\emph{Démonstration}. On peut supposer $Z$ et $X$ de dimension
\hyperref[XIII.sec1]{1}, le 
cas où $X$ est de dimension nulle étant triviale grâce à l'hypothèse
faite sur $k$. Alors $Z=X\cup Y$, où $Y$ est le sous-schéma fermé
adhérence de $Z-X$. Soit $V=X\cap Y$ le schéma intersection, qui est
un schéma discret \ie localement artinien. Le morphisme $X\amalg
Y\rightarrow Z$ est fini et surjectif, donc un morphisme de descente
strict pour la catégorie des revêtements étales (\SGA \hyperrefext[IX][IX.cor4.7]{4.7}). On a
évidemment
\begin{equation*}
\fprod{(X\amalg Y)}{(X\amalg Y)}{Z}=X\amalg Y\amalg V^{(1)}\amalg
V^{(2)}\quad (V^{(i)}=V){\quoi},\tag{$\ast$}
\end{equation*}
où $\pr_i$ envoie $V^{(i)}$ dans le composant $X$ de $X\amalg Y$ et
$V^{(j)}$ dans le composant $Y$ pour $j\neq i$.

Soit $X'/X$ un revêtement étale de degré $n$, et soit $Y''/Y$ un
revêtement de degré $n$ (par exemple le revêtement \emph{trivial} de
degré $n$), de sorte que $X'\amalg Y''$ est\pageoriginale un
revêtement de $X\amalg Y$ de degré $n$. Il suffit de trouver des
données de descente pour ce revêtement et pour le morphisme $X\amalg
Y\rightarrow Z$. En regardant $(\ast)$, on voit qu'une donnée de
descente pour ce morphisme équivaut simplement à un isomorphisme
$$
\fprod{X'}{V}{X}=V'\xrightarrow{\sim} V''=\fprod{Y''}{V}{Y}{\quoi}.
$$
Puisque $k$ est séparablement clos, il en est de même de $k(z_i)$, et
par suite chaque revêtement étale de $V$ est complètement décomposé,
d'où $V'\simeq V''$, ce qui donne le résultat.

\begin{enonce*}[remark]{Remarque 2.4}\etiquette[2.4]{XIII.rem2.4}
On peut éviter le recours au délicat résultat de \SGA IX et à la
théorie de la descente, par un argument direct, montrant que la
catégorie des revêtements étales de $Z$ est équivalente à la catégorie
des triples $(X',Y',\theta)$, où $X'(Y')$ est un revêtement étale de
$X(Y)$ et $\theta$ un $V$-isomorphisme $\fprod{X'}{V}{X}\simeq
\fprod{Y'}{V}{Y}$. Ce résultat, valable pour tout schéma $Z$ réunion
de deux sous-schémas fermés $X$, $Y$, s'établit directement à l'aide
du résultat analogue pour la détermination des Modules quasi-cohérents
sur $X$ (lorsque $X=V(\mathfrak{J})$, $Y=V(\bar{K})$, $J\cap \bar{K}=0$).
\end{enonce*}

\footnotetext[1]{Figure aussi dans \EGA IV 21.9.11, 21.9.12 sous une
forme légèrement plus générale; notamment il suffit dans \hyperref[XIII.cor3.2]{3.2} que $f$
soit \emph{séparé} au lieu de propre.}
\section{Un résultat auxiliaire sur le groupe de Picard ${}^{(*)}$}\etiquette[3]{XIII.sec3}

\begin{enonce*}{Proposition 3.1}\etiquette[3.1]{XIII.prop3.1}
Soient $X$ un schéma noethérien, $X_{\circ}$ un sous-schéma fermé de
$X$ tel que {\rm a)} $\dim X_{\circ}\leq 1$, et {\rm b)} pour toute
partie fermée connexe non-vide $Y$ de $X$, $Y_{\circ}=Y\cap X_{\circ}$
est connexe non-vide.
\begin{enumerate}
\renewcommand{\theenumi}{\roman{enumi}}
\renewcommand{\labelenumi}{(\theenumi)}
\item Soit\pageoriginale $Z_{\circ}$ la réunion des $\{\bar{x}\}\cap
X_{\circ}$ pour $x\in \Ass \Oo_X$ tel que $\{\bar{x}\}\cap X_{\circ}$
soit réduit à un point. Alors pour tout diviseur de Cartier
$D_{\circ}$ sur $X_{\circ}$ tel que $\supp$ ($D_{\circ}$ ne rencontre
pas $Z_{\circ}$, il existe un diviseur de Cartier $D$ sur $X$,
induisant $D_{\circ}$ (et si $D_{\circ}\geq 0$, on peut prendre $D\geq
0$).

\item Supposons qu'il existe un ouvert affine $U_{\circ}$ de
$X_{\circ}$ contenant l'ensemble fini $\Ass \Oo_{X_{\circ}}\cup
Z_{\circ}$ (condition automatiquement satisfaite $X$ est noethérien et
si $X_{\circ}$ admet un faisceau inversible ample). Alors
l'homomorphisme $\Pic(X)\rightarrow \Pic(X_{\circ})$ est surjectif.
\end{enumerate}
\end{enonce*}

\noindent
\emph{Démonstration}.
\begin{enumerate}
\renewcommand{\theenumi}{\roman{enumi}}
\renewcommand{\labelenumi}{(\theenumi)}
\item Comme $\dim X_{\circ}\leq 1$, on voit tout de suite que tout
diviseur de Cartier $D_{\circ}$ sur $X_{\circ}$ a un support discret,
donc est la différence de deux diviseurs de Cartier positifs, de
supports contenus dans celui de $D_{\circ}$. Cela nous ramène dans (i)
au cas où $D_{\circ}$ est positif, donc où pour tout $x\in \Supp
D_{\circ}$, $D_{\circ}$ est défini en $x$ par un élément non diviseur
de zéro $f_x\in \rad$ ($\Oo_{X_{\circ}},x$). Alors $f_x$ est induit
par une section $g_x$ de $\Oo_X$ sur un voisinage ouvert $U_x$ de
$x$. Je dis qu'en prenant $U_x$ assez petit, $g_x$ est non diviseur de
zéro dans $\Oo_{U_x}$, \ie, $V(g_x)\cap\Ass \Oo_{U_x}=∅$. En
effet, il résulte de a) et b) que pour tout $x\in X$, et en
particulier pour $z\in \Ass \Oo_X$, $\bar{z}\cap X_{\circ}$ est un
fermé connexe de $X_{\circ}$ donc est, soit réduit à un seul point
fermé, soit à une réunion connexe de composantes irréductibles de
dimension 1 sur $X_{\circ}$. Comme $x\not\in Z_{\circ}$ il s'ensuit
que dans le premier cas on a $x\not\in \bar{z}$. Mais pour $s\in \Ass
\Oo_{X}$ du deuxième type, on ne peut pas avoir $x\in \bar{z}$ et
$z\in V(g_x)$, puisqu'alors $g_x$ donc $f_x$ s'annulerait sur une des
composantes irréductibles de $\bar{z}\cap \supp X_{\circ}$, qui est
aussi une composante irréductible de $X_{\circ}$, ce qui est
incompatible avec le fait que $f_x$ est non diviseur de zéro dans
$\Oo_{X_{\circ},x}$. Donc quitte à restreindre $U_x$ on peut
supposer\pageoriginale que $V(g_x)\cap \Ass \Oo_{U_x}=\phi$, \ie
$g_x$ est non diviseur de zéro. Soit $Y_x$ l'adhérence de $V(g_x)$
dans $X$, alors $Y_x\cap X_{\circ}\cap U_x\subset V(g_x)\cap
X_{\circ}$ admet $x$ comme point isolé, donc il en est de même de
$Y_x\cap X_{\circ}$, donc par la condition b) de la proposition, $Y_x$
admet une composante connexe $Y'_x$ telle que $Y'_x\cap
X_{\circ}=\{x\}$. De plus, on a $Y'_x\subset U_x$, car
$T_x=Y'_n-Y'_x\cap U_x$ est un fermé de $X$ tel que $T_x\cap
X_{\circ}=\phi$, donc $T_x=\phi$ en vertu de la condition b). Nous
prenons maintenant le diviseur de Cartier $D'_x$ sur $X$ qui induit
$Y'_x$ sur $U_x$, et induit zéro sur $X-Y'_x$, et poserons $D=\Sigma_x
d'_x$, qui est un diviseur de Cartier positif sur $X$, induisant
évidemment $D_{\circ}$ sur $X_{\circ}$.

\item La conclusion résulte de (i) et du fait que tout module
inversible $\LL_{\circ}$ sur $X_{\circ}$ est isomorphe au module
défini par un diviseur de Cartier $D_{\circ}$ dont le support ne
rencontre pas $Z_{\circ}$. Ce dernier fait signifie en effet que pour
$\LL_{\circ}$ admet une section définie sur un ouvert $V_{\circ}$
contenant $\Ass \Oo_{X_{\circ}}\cup Z_{\circ}$, cette section ne
s'annulant en aucun point de ce dernier ensemble. Or il suffit de
prouver cela en remplaçant $X_{\circ}$ par l'ouvert affine
$U_{\circ}$, et dans ce cas cela est bien connu et provient du fait
qu'un module inversible sur un schéma semi-local est isomorphe au
faisceau structural. Enfin, on sait bien que si $X_{\circ}$ admet un
faisceau ample, alors toute partie finie de $X_{\circ}$ est contenue
dans un ouvert affine. Cela prouve (ii).
\end{enumerate}

On a le corollaire suivant, qui généralise \hyperrefext[XII][XII.prop8.6]{8.6} et qui achève la
démonstration de \hyperrefext[XII][XII.thm5.1]{5.1} (iii):

\begin{enonce*}{Corollaire 3.2}\etiquette[3.2]{XIII.cor3.2}
Soit $S$ le spectre d'un anneau hensélien noethérien, et
$f:X\rightarrow S$ un morphisme propre et de dimension relative $\leq
1$. Alors pour tout sous-schéma fermé $X'_{\circ}$ de $X$ ayant même
espace sous-jacent que la fibre fermée\pageoriginale $X_{\circ}$ de
$X/S$, l'application canonique $\Pic(X)\rightarrow \Pic(X_{\circ})$
est surjective.
\end{enonce*}

En effet, on applique la proposition à $(X,X'_{\circ})$, en notant que
la condition a) est satisfaite par hypothèse, et b) en vertu de 
\hyperrefext[XII][XII.cor5.7]{5.7}. Il reste à vérifier l'hypothèse faite dans ii)
d'existence d'un 
$U_{\circ}$, ce qui résulte du fait que $X'_{\circ}$ est
nécessairement projectif (comme toute courbe propre sur un corps).

\begin{enonce*}[remark]{Remarque}
Le lecteur qui ne voudra pas admettre ou reconstituer le fait que les
courbes propres sur un corps sont projectives, se bornera à utiliser
le corollaire en supposant déjà que $f$ est projectif, ce qui suffit
pour l'application que nous avions en vue ici. Signalons d'autre part
que, admettant la projectivité des courbes algébriques propres sur un
corps, on conclut que sous les conditions de \hyperref[XIII.cor3.2]{3.2} $X$ est
nécessairement projectif sur $S$, comme on voit en choisissant un
$\xi_{\circ}\in \Pic$ $(X_{\circ})$ ample et le relevant en un $\xi\in
\Pic(X)$, et appliquant \EGA III 4.7.1.
\end{enonce*}

\begin{thebibliography}{99}
\bibitem{SGA4-13:TDTE} Grothendieck A., Technique de descente et
théorèmes d'existence en Géométrie Algébrique, in Fondements de la
Géométrie Algébrique (Extraits du Sém. Bourbaki 1957-62) Secrétariat
mathématique, 11, rue Pierre Curie.
\end{thebibliography}

\ifx\danslelivre\undefined
\end{document}
\fi
